\label{chap:83-13}

\section{Construcci\'on preconstitutiva de la respuesta bidireccional de acci\'on y del operador de vac\'io}
\label{p12-83-c13-s1:ch:vacuum}

\subsection{Derivaci\'on de la escala de acci\'on a partir del estado dodecaf\'asico}

La d\'ecada de acci\'on se determina intr\'insecamente, sin adoptarla como
unidad heredada. El funcional determinantal local se aplica a la salida HMT
previamente construida \(\alpha_{\HMT}\) y define
\begin{equation}
\Sigma_H=\varphi_{\HMT}\alpha_{\HMT}^{16}.
\label{p12-83-c13-s1:eq:action-determinantal-reader}
\end{equation}
El exponente \(16\) es el orden del cokernel local del cierre firmado; su valoraci\'on decimal fija
\begin{equation}
\operatorname{ord}_{10}(\Sigma_H)=-34.
\label{p12-83-c13-s1:eq:action-decade}
\end{equation}

El origen del exponente se puede verificar dentro de este volumen. Sobre la
\'orbita local de cuatro hojas act\'ua el bloque de Hadamard
\begin{equation}
H_4=
\begin{pmatrix}
1&1&1&1\\
1&1&-1&-1\\
1&-1&1&-1\\
1&-1&-1&1
\end{pmatrix}.
\label{p12-83-c13-s1:eq:action-hadamard}
\end{equation}

\begin{lemma}[\'Indice local de la recta de acci\'on]
La forma normal de Smith de \(H_4\) es
\begin{equation}
\operatorname{SNF}(H_4)=\operatorname{diag}(1,2,2,4).
\label{p12-83-c13-s1:eq:action-smith}
\end{equation}
En consecuencia,
\[
\operatorname{coker}(H_4)\simeq
\Z/2\Z\oplus\Z/2\Z\oplus\Z/4\Z,
\qquad
|\operatorname{coker}(H_4)|=16.
\]
\end{lemma}

\begin{proof}
El m\'aximo com\'un divisor de las entradas es uno; el de los menores
\(2\times2\) no nulos es dos; el de los menores \(3\times3\) es cuatro; y
\(|\det H_4|=16\). Los cocientes sucesivos de estos divisores determinantes
son \(1,2,2,4\), que prueban \cref{p12-83-c13-s1:eq:action-smith}. El producto de los
invariantes da el orden del con\'ucleo.
\end{proof}

La norma de la secci\'on \(\alpha_{\HMT}\) sobre esas diecis\'eis hojas
produce \(\alpha_{\HMT}^{16}\); la autoescala \(\varphi_{\HMT}\) act\'ua
una sola vez sobre la base, no una vez por hoja. De ah\'i
\cref{p12-83-c13-s1:eq:action-determinantal-reader}. Las comparaciones certificadas
\begin{equation}
\alpha_{\HMT}^{16}<10^{-34}
<\varphi_{\HMT}\alpha_{\HMT}^{16}<10^{-33}
\label{p12-83-c13-s1:eq:action-decade-bounds}
\end{equation}
demuestran, sin introducir la unidad SI, que el orden decimal es
exactamente \(-34\). La evaluaci\'on terminal comienza por
\begin{equation}
\Sigma_H
=1.046243838104756580184229208303089\ldots\times10^{-34}.
\label{p12-83-c13-s1:eq:action-sigma-value}
\end{equation}
La sustituci\'on del con\'ucleo por tres copias con exponente \(16^3\), o la
aplicaci\'on reiterada de \(\varphi\) en cada hoja, modificar\'ia la
genealog\'ia. Ambas operaciones constituyen criterios de refutaci\'on y no
convenciones equivalentes.
La cadena causal que se conserva es, por tanto,
\[
\APP\to\TRIT\to\TPK\to U_{12}^{\rm sign}
\to\alpha_{\HMT}\to\Sigma_H
\to10^{-34}\U_S\to
\{\hbar_{\rm pre},\hbar_{\rm ret}\}.
\]
Este volumen no rederiva \(\alpha_{\HMT}\), sino que desarrolla la geometr\'ia
metrol\'ogica posterior a la selecci\'on del tipo y de la d\'ecada por
\(\Sigma_H\). El entero \(-34\) no es una mantisa ni un ajuste al Sistema
Internacional, sino la valoraci\'on terminal del funcional local.

\subsection{Secciones orientadas de la recta de acci\'on}

Las realizaciones anterior y posterior al retorno constituyen dos secciones
de una misma recta de acci\'on:

\begin{equation}
\hbar_{\rm pre}=H_5\,10^{-34}\U_S,\qquad
\hbar_{\rm ret}=\eta_{\rm ret}\,10^{-34}\U_S.
\label{p12-83-c13-s1:eq:action-twoway}
\end{equation}

La d\'ecada com\'un determina la recta, mientras que la informaci\'on
orientada se expresa mediante

\begin{equation}
R_{\rm act}=\frac{\hbar_{\rm pre}}{\hbar_{\rm ret}}
=\frac{H_5}{\eta_{\rm ret}},
\qquad
\xi_{\rm act}=\frac12\log R_{\rm act}.
\label{p12-83-c13-s1:eq:action-rapidity}
\end{equation}

La primera cantidad proporciona una coordenada multiplicativa y la segunda
una coordenada aditiva de tipo rapidez. No representan dos constantes de
Planck independientes, sino las dos orientaciones de una misma secci\'on
respecto del retorno.

Si \(h_a=2\pi\hbar_a\), \(a\in\{\mathrm{pre},\mathrm{ret}\}\), una magnitud proporcional a \(h_a^{1/2}\) cambia como \(R_{\rm act}^{1/2}\), una proporcional a \(h_a^{-1/2}\) cambia como \(R_{\rm act}^{-1/2}\), y un cociente donde la acci\'on se cancela es bidireccional invariante.

\subsection{Descomposici\'on espectral bicapa del operador angular}

Sea \(R=R^*=R^{-1}\) la involuci\'on activa y

\[
P_\pm=\frac{I\pm R}{2},
\qquad
T=q_+P_+ +q_-P_-,
\qquad 0<q_+<q_-<1.
\]

Con
\[
x=A_{\rm rad}=\frac{\pi A}{180},\qquad
y=C^*_{\rm rad}=\frac{\pi C^*}{180},
\]
los dos canales se construyen antes de evaluar el vac\'io:
\begin{equation}
q_+=e^{-x-y},\qquad q_-=e^{-x+y}.
\label{p12-83-c13-s1:eq:vacuum-angular-channels}
\end{equation}
El intercambio de hojas invierte \(C^*\), es decir \(y\mapsto-y\), y permuta \(q_+\leftrightarrow q_-\). La pareja es as\'i la lectura espectral de un mismo antecedente angular, no dos par\'ametros ajustados por separado.

\begin{definition}[Respuesta constitutiva]
\begin{equation}
\mathcal V_{90,120}=(I-T^{90})(I-T^{120})^{-1}.
\label{p12-83-c13-s1:eq:vacuum-operator}
\end{equation}
\end{definition}

Como \(\|T\|<1\), el denominador es positivo e invertible. El c\'alculo funcional da

\begin{equation}
\mathcal V_{90,120}=r_+P_+ +r_-P_-,
\qquad
r_\pm=\frac{1-q_\pm^{90}}{1-q_\pm^{120}}.
\label{p12-83-c13-s1:eq:vacuum-eigenvalues}
\end{equation}

La afirmaci\'on de invertibilidad merece hacerse expl\'icita. Para cada hoja,
\(0<q_\pm^{120}<1\); por tanto
\[
I-T^{120}=(1-q_+^{120})P_++(1-q_-^{120})P_-
\]
es estrictamente positivo y
\[
(I-T^{120})^{-1}
=\frac{P_+}{1-q_+^{120}}+
 \frac{P_-}{1-q_-^{120}}.
\]
La respuesta \(\mathcal V_{90,120}\) no se obtiene resolviendo cuatro
ecuaciones escalares separadas: es el resultado del c\'alculo funcional de una
sola contracci\'on positiva. Esta observaci\'on impide que permeabilidad,
permitividad, impedancia y propagaci\'on pierdan el antecedente que comparten.

\begin{theorem}[Determinaci\'on espectral conjunta de las cuatro relaciones constitutivas]
Los cuatro caracteres constitutivos
\begin{equation}
\widehat\varepsilon=r_+^2,\qquad
\widehat\mu=r_-^2,\qquad
\widehat Z=\frac{r_-}{r_+},\qquad
\widehat c=\frac1{r_+r_-}
\label{p12-83-c13-s1:eq:vacuum-four}
\end{equation}
son funciones espectrales de \(\mathcal V_{90,120}\). Satisfacen
\begin{equation}
\widehat\mu\widehat\varepsilon=\widehat c^{-2},
\qquad
\frac{\widehat\mu}{\widehat\varepsilon}=\widehat Z^2.
\label{p12-83-c13-s1:eq:vacuum-relations}
\end{equation}
\end{theorem}

\begin{proof}
Las identidades se obtienen sustituyendo \cref{p12-83-c13-s1:eq:vacuum-four}. No interviene ninguna carta dimensional ni ning\'un valor exterior.
\end{proof}

\begin{corollary}[Reconstrucci\'on de las dos hojas]
Los cuatro caracteres no contienen cuatro grados de libertad. Cualquiera de
los pares \((\widehat\mu,\widehat\varepsilon)\),
\((\widehat Z,\widehat c)\) o \((\mathcal E,\mathcal B)\) reconstruye los
dos autovalores positivos:
\begin{align}
r_-&=\sqrt{\widehat\mu}
=\sqrt{\widehat Z/\widehat c},
&r_+&=\sqrt{\widehat\varepsilon}
=\frac1{\sqrt{\widehat Z\widehat c}},
\label{p12-83-c13-s1:eq:vacuum-reconstruct-r}\\
\log r_-&=\frac{\mathcal E+\mathcal B}{2},
&\log r_+&=\frac{\mathcal E-\mathcal B}{2}.
\label{p12-83-c13-s1:eq:vacuum-reconstruct-log}
\end{align}
En particular, la orientaci\'on se pierde si se conserva s\'olo el producto,
pero se recupera al a\~nadir el cociente.
\end{corollary}

\begin{proof}
La primera l\'inea es la ra\'iz positiva de las identidades
\cref{p12-83-c13-s1:eq:vacuum-four,p12-83-c13-s1:eq:vacuum-relations}. La segunda invierte el sistema
lineal
\(\mathcal E=\log r_-+\log r_+\),
\(\mathcal B=\log r_--\log r_+\).
\end{proof}

Este corolario expresa una propiedad central de HMT: producto y cociente son
funcionales complementarios de un mismo estado. El producto determina la
propagaci\'on com\'un, mientras que el cociente conserva la orientaci\'on
relativa de las hojas interior y exterior.

En el certificado V2,

\begin{align*}
r_+&=0.9999996285482493631786952281\ldots,\\
r_-&=0.9997241371895183641315165853\ldots,\\
\widehat\varepsilon&=0.9999992570966367027604416155\ldots,\\
\widehat\mu&=0.9994483504793269350899815559\ldots,\\
\widehat Z&=0.9997245085389372154555543466\ldots,\\
\widehat c&=1.0002763104861575261063862689\ldots.
\end{align*}

Estos numerales son el reconocimiento terminal de \cref{p12-83-c13-s1:eq:vacuum-operator}; no definen sus exponentes \(90,120\).

\subsection{Factorizaci\'on arm\'onica de los sectores \texorpdfstring{\(90\)}{90} y \texorpdfstring{\(120\)}{120}}

El semigrupo arm\'onico es

\[
\langle90,120\rangle=30\langle3,4\rangle.
\]

Si \(s=q^{30}\), entonces

\begin{equation}
\frac{1-q^{90}}{1-q^{120}}
=\frac{1-s^3}{1-s^4}
=\frac{1+s+s^2}{1+s+s^2+s^3}.
\label{p12-83-c13-s1:eq:three-four-completion}
\end{equation}

La identidad tambi\'en separa la respuesta de su defecto. Si
\[
F_{3\to4}(s)=\frac{1+s+s^2}{1+s+s^2+s^3},
\qquad 0<s<1,
\]
entonces
\begin{equation}
d_{3\to4}(s):=1-F_{3\to4}(s)
=\frac{s^3}{1+s+s^2+s^3}>0.
\label{p12-83-c13-s1:eq:vacuum-defect34}
\end{equation}
El t\'ermino adicional se conserva exactamente como defecto
positivo. Adem\'as,
\begin{equation}
F_{3\to4}'(s)
=-\frac{s^2(3+2s+s^2)}{(1+s+s^2+s^3)^2}<0.
\label{p12-83-c13-s1:eq:vacuum-monotonic34}
\end{equation}
Por consiguiente, el orden de los canales angulares determina un\'ivocamente
el orden de las respuestas constitutivas. La asignaci\'on de
\(\widehat\mu\) y \(\widehat\varepsilon\) a las hojas respectivas se obtiene
anal\'iticamente, sin recurrir a sus expansiones decimales.

Cada hoja del vac\'io constituye la completaci\'on exacta de tres a cuatro
t\'erminos de un mismo modo elemental \(30\). En consecuencia, permeabilidad
y permitividad se interpretan como dos funcionales cuadr\'aticos de una misma
completaci\'on evaluada sobre hojas conjugadas, y no como par\'ametros
constitutivos independientes.

Como la funci\'on \((1-q^{90})/(1-q^{120})\) decrece en \(0<q<1\) y \(q_+<q_-\), se obtiene

\[
0<r_-<r_+<1,\qquad
0<\widehat\mu<\widehat\varepsilon<1,\qquad
0<\widehat Z<1<\widehat c.
\]

Los l\'imites del funcional determinan su geometr\'ia. Cuando \(q\to0^+\),
\(F(q)\to1\): las cuatro lecturas se acercan a la unidad y el defecto de
memoria desaparece. Cuando \(q\to1^-\),
\[
F(q)\longrightarrow\frac{90}{120}=\frac34.
\]
La respuesta toma, por tanto, valores en el intervalo positivo \((3/4,1)\). El
extremo \(3/4\) no es una constante ajustada: es el cociente de las dos
longitudes de canal y la realizaci\'on continua de la completaci\'on
tres--a--cuatro.

\subsection{Involuci\'on de hojas y conservaci\'on de la orientaci\'on}

El intercambio de hojas \(C^*\mapsto-C^*\) act\'ua por

\begin{equation}
\widehat\mu\longleftrightarrow\widehat\varepsilon,\qquad
\widehat Z\longmapsto\widehat Z^{-1},
\qquad
\widehat c\longmapsto\widehat c.
\label{p12-83-c13-s1:eq:vacuum-sheet}
\end{equation}

La propagaci\'on com\'un es par; la impedancia conserva la orientaci\'on.
Esta distinci\'on reaparecer\'a en Barbero y CKM: no todos los funcionales
asociados al cambio de hoja tienen la misma paridad.

\subsection{Caracteres logar\'itmicos y compatibilidad tracial}

Definamos

\[
\mathcal E=\log(r_+r_-),\qquad
\mathcal B=\log(r_-/r_+).
\]

Con la traza normalizada \(\tau=\Tr/2\),

\begin{equation}
2\tau(\log\mathcal V)=\mathcal E,
\qquad
-2\tau(R\log\mathcal V)=\mathcal B.
\label{p12-83-c13-s1:eq:vacuum-traces}
\end{equation}

La pareja \((\mathcal E,\mathcal B)\) constituye un sistema de coordenadas lineal del logaritmo
operatorio:
\begin{equation}
\log\mathcal V
=\frac{\mathcal E}{2}I-\frac{\mathcal B}{2}R.
\label{p12-83-c13-s1:eq:vacuum-log-decomposition}
\end{equation}
La componente escalar es invariante bajo el intercambio de hojas; la
componente orientada cambia de signo. Esta descomposici\'on demuestra
directamente la paridad de la propagaci\'on y el car\'acter orientado de la
impedancia.

La amplificaci\'on non\'adica

\[
\mathcal V_m=\mathcal V\otimes I_{9^m}
\]

es compatible con las expectativas \(E_m=\mathrm{id}\otimes\tau_9\):

\[
E_m(\mathcal V_{m+1})=\mathcal V_m.
\]

Por tanto, \(\mathcal E\) y \(\mathcal B\) no son peculiaridades de una matriz \(2\times2\); permanecen en la torre UHF y en su cierre tracial \(II_1\).

M\'as precisamente, para todo polinomio \(p\) y despu\'es, por continuidad,
para toda funci\'on continua en el espectro,
\begin{equation}
E_m\bigl(p(\mathcal V_{m+1})\bigr)=p(\mathcal V_m).
\label{p12-83-c13-s1:eq:vacuum-functional-refinement}
\end{equation}
La compatibilidad no preserva solamente dos n\'umeros: preserva todo el
\'algebra funcional generada por la respuesta. En particular conserva sus
proyectores espectrales, sus potencias, su logaritmo y las formas cuadr\'aticas
que miden los defectos \cref{p12-83-c13-s1:eq:vacuum-defect34}.

\subsection{Realizaci\'on dimensional de las relaciones constitutivas}

La Mec\'anica Dimensional determina

\begin{align}
Z_{\rm vac}&=\widehat Z\,u_Su_Q^{-2},\\
c_{\rm vac}&=\widehat c\,u_C,\\
\mu_{\rm vac}&=\widehat\mu\,u_Su_C^{-1}u_Q^{-2},\\
\varepsilon_{\rm vac}&=\widehat\varepsilon\,u_S^{-1}u_C^{-1}u_Q^2.
\label{p12-83-c13-s1:eq:vacuum-sections}
\end{align}

Las identidades constitutivas se mantienen en las rectas:

\begin{equation}
\mu_{\rm vac}\varepsilon_{\rm vac}=c_{\rm vac}^{-2},
\qquad
\mu_{\rm vac}/\varepsilon_{\rm vac}=Z_{\rm vac}^2.
\end{equation}

La carta SI es posterior. No es necesario declarar \(c\) primitivo y resolver despu\'es \(\mu\varepsilon=c^{-2}\): los tres objetos descienden conjuntamente de \(\mathcal V\) y de las rectas declaradas.

La separaci\'on de niveles se expresa mediante el diagrama
\[
\begin{array}{c}
(A,C^*)\longrightarrow T\longrightarrow\mathcal V
\longrightarrow
(\widehat\mu,\widehat\varepsilon,\widehat Z,\widehat c)
\\[2mm]
\Big\downarrow\text{ cartas de recta}\hspace{36mm}
\Big\downarrow\text{ evaluaci\'on final}
\\[2mm]
(u_S,u_Q,u_C)\longrightarrow
(\mu_{\rm vac},\varepsilon_{\rm vac},Z_{\rm vac},c_{\rm vac})
\longrightarrow\text{ coordenadas SI}.
\end{array}
\]
La primera l\'inea es adimensional y operatoria; la segunda determina tipos de
magnitud; la tercera selecciona una carta metrol\'ogica. Alterar la carta final
no modifica ni los autovalores ni las identidades constitutivas. En ese
sentido preciso, la unificaci\'on matem\'atica--f\'isica consiste en distinguir
el invariante geneal\'ogico de su sistema de coordenadas, y no en suprimir la
estructura dimensional.

\begin{proposition}[Minimalidad espectral]
Sea \(W\) un operador positivo de dos hojas con espectro simple
\(\{r_+,r_-\}\). Si cuatro caracteres positivos satisfacen
\(\mu\varepsilon=c^{-2}\) y \(\mu/\varepsilon=Z^2\), si la propagaci\'on
com\'un se determina mediante \(c^{-1}=\det W=r_-r_+\), y si su orientaci\'on
fija \(Z=r_-/r_+\), entonces necesariamente
\[
(\mu,\varepsilon,Z,c)
=\left(r_-^2,r_+^2,\frac{r_-}{r_+},\frac1{r_-r_+}\right).
\]
Por tanto \cref{p12-83-c13-s1:eq:vacuum-four} no es una parametrizaci\'on entre muchas:
es la realizaci\'on positiva m\'inima compatible con producto, cociente y
orientaci\'on.
\end{proposition}

\begin{proof}
De \(Z=r_-/r_+\) y \(c^{-1}=\det W=r_-r_+\) se obtienen de forma \`unica las
ra\'ices positivas. Las dos identidades constitutivas fuerzan despu\'es
\(\mu=r_-^2\) y \(\varepsilon=r_+^2\).
\end{proof}

\subsection{Determinante angular y relaci\'on electrod\'ebil de Weinberg}

El determinante de \(T\) es

\begin{equation}
\det T=q_+q_-=D_A.
\label{p12-83-c13-s1:eq:detT}
\end{equation}

El sector constitutivo del vac\'io eval\'ua \(F(T)^2\), donde
\(F(z)=(1-z^{90})/(1-z^{120})\), mientras que el sector angular
electrod\'ebil eval\'ua \(\det T\). Se trata de dos funciones de un
antecedente com\'un y no de una factorizaci\'on horizontal entre constantes.
Esta distinci\'on se desarrolla en el \cref{ch:ew}.

\begin{falsifier}
El bloque queda refutado si se viola cualquiera de estas identidades exactas:
positividad de \(I-T^{120}\), igualdad
\cref{p12-83-c13-s1:eq:three-four-completion}, relaciones \cref{p12-83-c13-s1:eq:vacuum-relations} o
compatibilidad \(E_m(\mathcal V_{m+1})=\mathcal V_m\). Una coincidencia SI
no puede compensar un fallo operatorio.
\end{falsifier}

\begin{statusbox}
Operador, espectro, completaci\'on \(3\to4\), paridad, amplificaci\'on y
secciones constitutivas forman un teorema interno exacto. Su reuni\'on en un
solo operador queda establecida en esta formulaci\'on; los datos angulares y
las rectas, derivados previamente en HMT, constituyen la estructura de partida
expl\'icita.
\end{statusbox}


\section{Constantes cu\'antico--el\'ectricas, escala t\'ermica y sistema de Planck}
\label{p12-83-c13-s2:ch:metrology}

\subsection{Derivaci\'on de la secci\'on de carga a partir del vac\'io}

Fijados \(Z_{\rm vac}\), \(h_a=2\pi\hbar_a\) y \(\alpha_{\HMT}\), la carga de la hoja \(a\in\{\mathrm{pre},\mathrm{ret}\}\) es

\begin{equation}
e_a=\sqrt{\frac{2\alpha_{\HMT}h_a}{Z_{\rm vac}}}.
\label{p12-83-c13-s2:eq:charge}
\end{equation}

La secci\'on \(e_a\) no se introduce mediante una coordenada del Sistema
Internacional destinada a reconstruir \(\alpha\). El orden causal procede de
\(\alpha_{\HMT}\), de la acci\'on y del vac\'io hacia la secci\'on de carga.

\subsection{Constantes de von Klitzing y Josephson, conductancia y flujo magn\'etico}

De \cref{p12-83-c13-s2:eq:charge} se deducen

\begin{align}
R_K&=\frac{h_a}{e_a^2}=\frac{Z_{\rm vac}}{2\alpha_{\HMT}},
\label{p12-83-c13-s2:eq:rk}\\
G_0&=\frac{2e_a^2}{h_a}=\frac{4\alpha_{\HMT}}{Z_{\rm vac}},
\label{p12-83-c13-s2:eq:g0}\\
\Phi_{0,a}&=\frac{h_a}{2e_a}
=\sqrt{\frac{h_aZ_{\rm vac}}{8\alpha_{\HMT}}},
\label{p12-83-c13-s2:eq:phi0}\\
K_{J,a}&=\frac{2e_a}{h_a}=\Phi_{0,a}^{-1}.
\label{p12-83-c13-s2:eq:kj}
\end{align}

En consecuencia,

\begin{equation}
R_KG_0=2,\qquad G_0Z_{\rm vac}=4\alpha_{\HMT},
\qquad K_{J,a}\Phi_{0,a}=1.
\label{p12-83-c13-s2:eq:quantum-electrical-identities}
\end{equation}

La covariancia bidireccional adopta la forma

\[
R_K^{\rm pre}=R_K^{\rm ret},\qquad
G_0^{\rm pre}=G_0^{\rm ret},
\]

\[
\frac{\Phi_{0,\rm pre}}{\Phi_{0,\rm ret}}=R_{\rm act}^{1/2},
\qquad
\frac{K_{J,\rm pre}}{K_{J,\rm ret}}=R_{\rm act}^{-1/2}.
\]

Las relaciones de resistencia y conductancia son independientes de la
orientaci\'on de la acci\'on. El par flujo--frecuencia conserva dicha
orientaci\'on y transforma rec\'iprocamente sus dos componentes.

\subsection{Relaci\'on t\'ermica intr\'inseca de la escala \texorpdfstring{\(108\)}{108}}

La estructura cronol\'ogica \(108\) fija

\begin{equation}
k_B\Theta_{\rm clk}\log3
=\frac{h_{\rm int}}{108t_0}
=\frac{2\pi\hbar_{\rm int}}{108t_0}.
\label{p12-83-c13-s2:eq:boltzmann-clock}
\end{equation}

Esta igualdad no determina por separado una coordenada num\'erica de \(k_B\) y otra de \(\Theta_{\rm clk}\) sin una normalizaci\'on adicional. S\'i determina su producto como secci\'on de energ\'ia. El factor \(\log3\) tiene tres generadors compatibles:

\begin{equation}
\log3=C_1+C_2-2C_0
=\frac{S_{\TRIT}}{k_B}
=\frac{E_P}{k_B\Theta_{\rm clk}}.
\label{p12-83-c13-s2:eq:log3-triple}
\end{equation}

Las tres igualdades corresponden a realizaciones residual, informacional y
termodin\'amica de un mismo escalar. La coincidencia del valor no identifica
sus dominios respectivos.

\subsection{Sistema de Planck y covariancia bidireccional}

El selector CT108 construye

\begin{equation}
\theta_{108}=\left(\frac{54}{\pi}\right)^2,
\qquad
\ell_P=\frac{54}{\pi}\ell_0,
\qquad
t_P=\frac{54}{\pi}t_0,
\label{p12-83-c13-s2:eq:planck-length-time}
\end{equation}

y

\begin{equation}
E_P=\frac{\hbar_{\rm int}}{t_P}
=\frac{\pi\hbar_{\rm int}}{54t_0}.
\label{p12-83-c13-s2:eq:planck-energy}
\end{equation}

Comparando \cref{p12-83-c13-s2:eq:boltzmann-clock,p12-83-c13-s2:eq:planck-energy} se obtiene

\begin{equation}
\boxed{E_P=k_B\Theta_{\rm clk}\log3.}
\label{p12-83-c13-s2:eq:planck-trit}
\end{equation}

La relaci\'on no depende de una coincidencia entre unidades: ambos miembros
proceden de la misma estructura cronol\'ogica \(108\).

Cuando la constante gravitatoria del \cref{ch:gravity} se introduce, la familia de Planck satisface

\begin{align}
\ell_P^2&=\frac{G\hbar}{c^3}=\theta_{108}\ell_0^2,
\label{p12-83-c13-s2:eq:planck-area}\\
F_P&=\frac{c^4}{G},
\qquad
P_P=\frac{c^5}{G},
\label{p12-83-c13-s2:eq:planck-force-power}\\
\rho_P&=\frac{\hbar}{\theta_{108}^2c^5t_0^4}.
\label{p12-83-c13-s2:eq:planck-density}
\end{align}

Aunque \(G\) y \(\hbar\) se transformen rec\'iprocamente entre las secciones
anterior y posterior al retorno, su producto en \(\ell_P^2\) conserva la
geometr\'ia. Esta invariancia determina el \`area como interfaz entre las
realizaciones geom\'etrica, informacional y de entrelazamiento.

\subsection{Formulaciones espectrales de la ley de desplazamiento de Wien}

Para la densidad espectral por longitud de onda, el m\'aximo satisface

\[
5(1-\ee^{-x_5})=x_5,
\qquad
x_5=5+W_0(-5\ee^{-5}).
\]

Con \cref{p12-83-c13-s2:eq:boltzmann-clock},

\begin{equation}
\lambda_{\max}(\Theta_{\rm clk})
=\frac{108\log3}{x_5}\,\ell_0.
\label{p12-83-c13-s2:eq:wien-lambda}
\end{equation}

Para la densidad por frecuencia,

\[
3(1-\ee^{-x_3})=x_3,
\qquad
x_3=3+W_0(-3\ee^{-3}),
\]

y

\begin{equation}
\nu_{\max}(\Theta_{\rm clk})
=\frac{x_3}{108t_0\log3}.
\label{p12-83-c13-s2:eq:wien-frequency}
\end{equation}

Los m\'aximos no son rec\'iprocos, puesto que las densidades espectrales se
transforman con el jacobiano correspondiente. El cociente \(x_5/x_3\)
pertenece a la estructura de la parametrizaci\'on y no representa una
inconsistencia.

\subsection{Ley de Stefan--Boltzmann y termodin\'amica del gas de fotones}

Evaluar la ley de radiaci\'on en \(\Theta_{\rm clk}\) produce el flujo intr\'inseco

\begin{equation}
j_\star(\Theta_{\rm clk})
=\frac{2\pi^5}{15\,108^4(\log3)^4}\,
\frac{h_{\rm int}}{\ell_0^2t_0^2},
\label{p12-83-c13-s2:eq:stefan-hmt}
\end{equation}

en la normalizaci\'on declarada. La constante de Stefan--Boltzmann no se inserta como primitiva: resulta de integrar el espectro bos\'onico una vez fijados acci\'on, propagaci\'on y temperatura.
Las leyes espectrales empleadas son las de Planck y Wien \cite{rm:Planck,rm:Wien}; HMT fija aqu\'i la secci\'on sobre la que se eval\'uan, no sus constantes por contraste retrospectivo.

En unidades de la escala de Planck,

\begin{align}
n_\gamma\ell_P^3
&=\frac{2\zeta(3)}{\pi^2\log^3 3},
\label{p12-83-c13-s2:eq:photon-number}\\
\frac{u_\gamma\ell_P^3}{E_P}
&=\frac{\pi^2}{15\log^4 3},
\label{p12-83-c13-s2:eq:photon-energy}\\
\frac{s_\gamma\ell_P^3}{k_B}
&=\frac{4\pi^2}{45\log^3 3}.
\label{p12-83-c13-s2:eq:photon-entropy}
\end{align}

\(\zeta(3)\) representa el momento bos\'onico de orden tres, mientras que
\(\log3\) determina la energ\'ia de decisi\'on TRIT. La ecuaci\'on conserva la
distinci\'on entre ambos invariantes.

La conductancia t\'ermica cu\'antica se expresa como

\begin{equation}
\frac{\kappa_Qt_0}{k_B}=\frac{\pi^2}{324\log3}.
\label{p12-83-c13-s2:eq:thermal-conductance}
\end{equation}

\subsection{Confluencia metrológica posterior de acción, carga, conductancia y flujo magnético}

\begin{longtable}{@{}p{0.18\textwidth}p{0.27\textwidth}p{0.45\textwidth}@{}}
\toprule
Objeto & Ecuaci\'on interna & Significado estructural\\
\midrule
\endhead
\(R_K\) & \(Z_{\rm vac}/(2\alpha)\) & resistencia cu\'antica bidireccional invariante\\
\(G_0\) & \(4\alpha/Z_{\rm vac}\) & conductancia complementaria\\
\(\Phi_0\) & \(\sqrt{hZ/(8\alpha)}\) & flujo que conserva orientaci\'on de acci\'on\\
\(K_J\) & \(\Phi_0^{-1}\) & relaci\'on rec\'iproca entre frecuencia y flujo\\
\(E_P\) & \(k_B\Theta\log3\) & energ\'ia del ciclo y contenido informacional TRIT\\
Wien & \(108\log3/x_5\) & extremo espectral en la escala cronol\'ogica\\
Ap\'ery fot\'onica & \(2\zeta(3)/(\pi^2\log^3 3)\) & densidad de n\'umero bos\'onica\\
\bottomrule
\end{longtable}

\begin{falsifier}
El bloque queda refutado si un cambio de variable espectral conserva
indebidamente el mismo extremo de Wien, si \(\zeta(3)\) sustituye a
\(\log3\), o si se asignan valores separados a \(k_B\) y \(\Theta\) sin
especificar una carta. Las identidades
\cref{p12-83-c13-s2:eq:quantum-electrical-identities,p12-83-c13-s2:eq:planck-trit} proporcionan controles
algebraicos directos.
\end{falsifier}

\begin{statusbox}
Cuantos el\'ectricos y confluencia Planck--TRIT: \status{Teorema interno exacto}. Wien, Stefan--Boltzmann y gas fot\'onico: \status{Composici\'on matem\'atica exacta} con el continuo bos\'onico declarado. La carta SI permanece como \status{Contraste metrol\'ogico externo}.
\end{statusbox}


\section{Realización analítica graduada de la microfibra quíntuple}
\label{p12-83-c13-s3:chap:realizacion-analitica-contraangulo}
\index{coordenada angular conjugada \(C^*\)}
\index{microfibra de pi@microfibra de \(\pi\)!carácter analítico}
\index{carácter determinantal}

El catálogo finito y una evaluación analítica responden a preguntas de tipos
distintos. El catálogo determina qué regiones existen, cuáles comparten un
retorno y qué información conserva cada una. Una evaluación analítica debe
decidir, además, cómo se transforma la longitud de una palabra en grado, cómo
se prolonga una rama después de su retorno y con qué carácter escalar se lee
esa prolongación. Este capítulo construye esa segunda operación de manera
explícita.

La construcción parte de cuatro datos ya obtenidos en los capítulos
anteriores: la microfibra de cinco regiones de \(\pi\), la longitud seis de
cada palabra regional, el cierre dodecafásico de \(\alpha\) y el transporte
bicapa asociado a la coordenada común
\(A=10^3\alpha\). Ningún valor de la coordenada angular conjugada \(C_*^{\rm HMT}\), de la componente reducida de
acción, del funcional de Barbero--Immirzi o de una constante metrológica se
emplea en la evaluación. El resultado es un carácter regional convergente,
una recurrencia exacta y una fase orientada cuya carta sexagesimal reproduce
la coordenada angular conjugada publicada.

La separación entre la parte finita y la parte analítica es esencial. El
entero \(169\) será un teorema del catálogo. La serie que lo prolonga será el
teorema de un lector analítico cuyas reglas se declaran antes de evaluarlo.
Así, las premisas adicionales no quedan ocultas dentro de una coincidencia
decimal.

\subsection{El retorno persistente de las \(5\) regiones}

Sea
\[
\mathcal F_\pi=\Psi^{-1}(010211)\subset\mathcal U_6
\]
la microfibra relativa al catálogo. Recordemos sus cinco elementos, ahora
separados en bloque inicial y bloque terminal:
\[
\begin{aligned}
 &(501,614,498\mid169,272,272),\\
 &(810,923,870\mid169,272,272),\\
 &(810,983,810\mid169,272,272),\\
 &(870,923,810\mid169,272,272),\\
 &(870,923,810\mid418,674,521).
\end{aligned}
\]
Definimos la proyección de retorno
\[
p_{\rm ret}:\mathbb Z^6\longrightarrow\mathbb Z^3,
\qquad
p_{\rm ret}(u_1,\ldots,u_6)=(u_4,u_5,u_6).
\]

\begin{proposition}[Bloque terminal persistente]
\label{p12-83-c13-s3:prop:bloque-terminal-persistente}
El multiconjunto \(p_{\rm ret}(\mathcal F_\pi)\) posee un único elemento de
multiplicidad máxima:
\[
b_\pi=(169,272,272),
\qquad
\operatorname{mult}_{\mathcal F_\pi}(b_\pi)=4.
\]
El único bloque terminal restante es
\[
b_\pi^-=(418,674,521)
\]
y tiene multiplicidad uno.
\end{proposition}

\begin{proof}
La afirmación se obtiene por enumeración exacta de las cinco filas del
catálogo que satisfacen \(\Psi(U)=010211\). Cuatro filas poseen bloque
terminal \((169,272,272)\); la quinta, que porta el retorno mutado, posee
\((418,674,521)\). La multiplicidad máxima es cuatro y se alcanza una sola
vez.
\end{proof}

La diferencia orientada entre el retorno mutado y el persistente es
\[
 \omega_\pi=b_\pi^- - b_\pi=(249,402,249),
 \qquad
 \langle(1,1,1),\omega_\pi\rangle=900.
\]
Esta identidad separa dos objetos que no deben confundirse. El entero
\(169\) es una coordenada de un dato terminal, es decir, un dato de grado
cero. El vector \(\omega_\pi\) puede leerse como una cocadena de grado uno
después de orientar la arista que une ambos bloques. Llamarlo cociclo exigiría,
además, fijar un complejo de cocadenas y verificar que su coborde se anula.
La prolongación analítica construida más adelante no presupone esa
identificación.

La selección de \(169\) no exige privilegiar la primera posición del bloque.
El grupo simétrico \(\mathfrak S_3\) actúa permutando sus tres coordenadas.
El estabilizador de \(b_\pi\) intercambia las dos apariciones de \(272\) y
fija la única coordenada de multiplicidad uno. Denotemos por
\(\operatorname{sing}(b)\) esa coordenada cuando el bloque tiene tipo
\((r,s,s)\), con \(r\ne s\). Equivalentemente,
\[
\operatorname{sing}(b)
=\sum_{j=1}^3b_j-2\operatorname{moda}(b).
\]
En particular,
\[
\boxed{\operatorname{sing}(b_\pi)=169.}
\]

\begin{definition}[Selector residual de la microfibra]
\label{p12-83-c13-s3:def:selector-residual-pi}
El selector \(\rho_\pi\) aplica sucesivamente dos operaciones:
\begin{enumerate}
  \item selecciona el bloque terminal de multiplicidad máxima dentro de
  \(p_{\rm ret}(\mathcal F_\pi)\);
  \item retiene la coordenada fija por el estabilizador de ese bloque.
\end{enumerate}
Por la \cref{p12-83-c13-s3:prop:bloque-terminal-persistente},
\[
\boxed{\rho_\pi=169.}
\]
\end{definition}

\begin{theorem}[Unicidad equivariante del retorno]
\label{p12-83-c13-s3:thm:unicidad-retorno-169}
Entre los selectores que
\begin{enumerate}
  \item son invariantes bajo permutaciones de las cinco regiones;
  \item son equivariantes bajo permutaciones de las tres coordenadas
  terminales;
  \item seleccionan el bloque de multiplicidad máxima;
  \item retienen una coordenada fija por su estabilizador,
\end{enumerate}
el valor de retorno está determinado de manera única y es \(169\).
\end{theorem}

\begin{proof}
La invariancia bajo \(\mathfrak S_5\) impide usar el orden de las filas. La
tercera condición selecciona \(b_\pi\), que es el único modo por la
\cref{p12-83-c13-s3:prop:bloque-terminal-persistente}. Su estabilizador en
\(\mathfrak S_3\) posee dos órbitas sobre las coordenadas: la pareja de
valores \(272\) y el valor unipuntual \(169\). La cuarta condición selecciona la
segunda órbita. La equivarianza conserva su valor al desplazar su posición.
\end{proof}

Este argumento mejora la lectura posicional del retorno. El \(169\) no es
«la primera coordenada» de una fila escogida: es un invariante del
multiconjunto regional y de la acción de sus simetrías de reordenación.
Tampoco se lo selecciona por rareza global. En las \(468\) emisiones del
catálogo, \(169\) aparece \(84\) veces, uniformemente distribuido con
\(14\) apariciones en cada una de las seis posiciones. La unicidad del
\cref{p12-83-c13-s3:thm:unicidad-retorno-169} es, por tanto, relacional y regional: reside
en la combinación de microfibra, persistencia modal y estabilizador, no en la
frecuencia aislada del entero.

\subsection{Longitud cilíndrica y grado analítico}

Sea \(\mathcal P\) el módulo libre generado por las palabras regionales y
sea \(F^n\mathcal P\) su filtración por longitud. La graduación asociada
registra el primer nivel en el que aparece cada palabra. Para un parámetro
\(z\), el carácter ordinario de longitud es
\[
\chi_z([w])=z^{|w|}.
\]
La multiplicatividad
\[
\chi_z([uv])=\chi_z([u])\chi_z([v])
\]
expresa que la concatenación suma longitudes. Al evaluar en
\(z=\alpha\), una palabra de longitud \(n\) recibe grado \(\alpha^n\).

\begin{definition}[Compatibilidad de grado]
\label{p12-83-c13-s3:def:compatibilidad-grado}
El lector analítico regional es compatible con la filtración cilíndrica
cuando envía el componente homogéneo de longitud \(n\) al grado analítico
\(n\):
\[
\operatorname{gr}_n\mathcal P\longrightarrow
\alpha^n\mathbb R.
\]
\end{definition}

Las cinco regiones pertenecen a \(\mathcal U_6\); por ello, el dato de
retorno aparece en grado seis:
\[
\rho_\pi[U_6]\longmapsto169\alpha^6.
\]
Este seis es la longitud de la palabra regional. No es la longitud de los
paseos cerrados del \cref{chap:cinco-ciclos-pi}. Aquellos paseos recorren a
la vez los cuatro anclajes
\(U_\pi,U_e,U_\varphi,U_{\rm APP}\) y poseen longitud mínima ocho. Una
palabra, una arista del grafo de bloques y un paseo cerrado son objetos de
tipos distintos; la graduación utilizada aquí pertenece al primero.

\subsection{El carácter determinantal del transporte común}

El cierre dodecafásico proporciona \(\alpha\). La completación de las tres
parejas nonádicas proporciona la escala \(10^3\), de modo que
\[
A=10^3\alpha.
\]
Elegida la carta angular arquimediana, la coordenada común en radianes es
\[
x=\frac{\pi A}{180}=\frac{50\pi}{9}\alpha.
\]
Sea \(R\) una involución real de traza nula en dimensión dos,
\[
R^2=I_2,
\qquad
\operatorname{tr}R=0,
\]
y consideremos el transporte formal bicapa
\[
\mathcal T(x,y)=\exp(-xI_2-yR).
\]
La variable orientada \(y\) todavía no se ha evaluado. Sin embargo, la acción
de \(\mathcal T\) sobre la línea determinante es independiente de ella:
\[
\begin{aligned}
\det\mathcal T(x,y)
&=\exp\!\left(\operatorname{tr}(-xI_2-yR)\right)\\
&=e^{-2x}.
\end{aligned}
\]
Definimos
\[
\boxed{
D_A=e^{-2x}=e^{-100\pi\alpha/9}.}
\]
El subíndice recuerda que se trata del determinante de la contracción común
asociada a \(A\), no de una función de \(C_*^{\rm HMT}\).

\begin{proposition}[Carácter de la línea determinante]
\label{p12-83-c13-s3:prop:caracter-determinante}
La aplicación
\[
\delta_A:(\mathbb N_0,+)\longrightarrow(\mathbb R_{>0},\cdot),
\qquad
\delta_A(k)=D_A^k,
\]
es el único carácter multiplicativo cuyo valor sobre una prolongación
elemental es \(D_A\).
\end{proposition}

\begin{proof}
La multiplicatividad da
\[
\delta_A(k)
=\delta_A(\underbrace{1+\cdots+1}_{k\text{ veces}})
=\delta_A(1)^k=D_A^k.
\]
La misma identidad prueba existencia y unicidad.
\end{proof}

El determinante permanece invariante bajo conjugación de
\(\mathcal T\) y bajo el intercambio de hojas \(R\mapsto-R\). Por tanto, la
cola construida a partir de \(D_A\) pertenece al modo común. La orientación
aparecerá en el signo con el que el carácter regional corrige la fase.

\subsection{Reglas del lector analítico regional}

La extensión analítica se fija mediante las reglas siguientes.

\begin{definition}[Lector regional determinantal]
\label{p12-83-c13-s3:def:lector-regional-determinantal}
El lector regional determinantal satisface:
\begin{enumerate}
  \item \emph{compatibilidad de grado}: la longitud cilíndrica es el grado
  analítico;
  \item \emph{descomposición de renovación}: el impulso finito de retorno y
  la continuación estricta pertenecen a sumandos aditivos distintos;
  \item \emph{normalización de continuación}: después de registrar una sola
  vez el residuo \(\rho_\pi\), la rama superviviente comienza en la unidad de
  la línea determinante;
  \item \emph{covariancia por concatenación}: cada prolongación elemental
  actúa mediante \(\bigwedge^2\mathcal T=D_A\);
  \item \emph{orientación}: en la convención cardinal fijada para APP, el
  retorno regional pertenece al sector compensador y se sustrae de la fase
  de acción de quinto grado.
\end{enumerate}
\end{definition}

La tercera regla tiene contenido matemático. El valor \(169\) se registra
una vez como amplitud del retorno; no se multiplica de nuevo en cada
prolongación. La continuación estricta cuenta la única línea que prosigue y
la normaliza por su unidad. Esta es la diferencia entre una recurrencia de
renovación y la iteración homogénea del impulso inicial.

La necesidad de declarar esa normalización puede probarse. Para todo
\(\lambda\in\mathbb R\), la familia
\[
F_\lambda(z)
=169z^6+\lambda\frac{D_Az^7}{1-D_Az}
\]
conserva el mismo retorno finito y la misma razón determinantal entre
coeficientes consecutivos. El catálogo y la razón por sí solos no distinguen
\(\lambda\). La unidad de la línea determinante fija
\(\lambda=1\) antes de evaluar \(z=\alpha\). De este modo, la normalización
no se obtiene ajustando el valor final; forma parte de la definición del
carácter.

\subsection{Recurrencia de retorno y suma cerrada}

Definimos los coeficientes \(\kappa_n\) por
\[
\kappa_6=169,
\qquad
\kappa_7=D_A,
\qquad
\kappa_{n+1}=D_A\kappa_n\quad(n\ge7).
\]
La primera igualdad es el impulso regional; las dos siguientes describen la
continuación normalizada.

\begin{theorem}[Realización analítica graduada]
\label{p12-83-c13-s3:thm:realizacion-analitica-graduada}
El germen analítico compatible con las reglas de la
\cref{p12-83-c13-s3:def:lector-regional-determinantal} es único y vale
\[
\boxed{
\mathscr C_\pi(z)
=169z^6+\frac{D_Az^7}{1-D_Az}.}
\]
Sus coeficientes satisfacen
\[
\kappa_n=D_A^{n-6}\qquad(n\ge7).
\]
\end{theorem}

\begin{proof}
Escribamos
\[
F(z)=169z^6+\sum_{k\ge1}c_kz^{6+k}.
\]
La normalización de continuación y la covariancia determinantal dan
\[
c_1=D_A,
\qquad
c_{k+1}=D_Ac_k.
\]
Por inducción, \(c_k=D_A^k\). La suma geométrica produce
\[
\sum_{k\ge1}D_A^kz^{6+k}
=\frac{D_Az^7}{1-D_Az}.
\]
No queda ningún coeficiente libre.
\end{proof}

Como \(0<\alpha<1\) y \(0<D_A<1\),
\[
0<D_A\alpha<1.
\]
En consecuencia, \(\mathscr C_\pi(\alpha)\) converge absolutamente y su
radio de convergencia es \(D_A^{-1}>1\). Numéricamente,
\[
D_A=0.7751291192471564301888840964802\ldots,
\]
\[
D_A\alpha=0.00565639046986492673885079095469\ldots.
\]
La rapidez de esta contracción explica por qué la suma completa y su
truncamiento de grado siete poseen las mismas quince primeras cifras en la
carta final.

\subsection{El carácter de acción de 5.º grado}

El bloque central
\[
\mathcal B_c=\begin{pmatrix}7&2\\2&7\end{pmatrix}
\]
posee autovalores \(9\) y \(5\). El ciclo dodecafásico aporta el rango
\(12\); sus órbitas de paso tres poseen longitud cuatro; y el censo completo
aporta el cociente exacto
\[
\frac{104976}{1944}=54.
\]
A partir de estos invariantes se fija el carácter de acción de quinto grado
\[
\boxed{
H_5(\alpha,\varphi)
=\frac12\sqrt{\frac{1000\alpha}{\varphi}}-\alpha
+\frac9{16}\alpha^2-\frac59\alpha^3
+\frac7{48}\alpha^4-\frac1{54}\alpha^5.}
\]
Los cuatro coeficientes racionales pueden escribirse como
\[
\frac9{16}=\frac{\lambda_+}{4^2},
\qquad
\frac59=\frac{\lambda_-}{\lambda_+},
\qquad
\frac7{48}=\frac{\tfrac12\operatorname{tr}\mathcal B_c}{4\cdot12},
\qquad
\frac1{54}=\frac{1944}{104976},
\]
con \((\lambda_+,\lambda_-)=(9,5)\). Estas identidades certifican la
procedencia de los números utilizados. La asignación sucesiva de esos
invariantes a los grados dos, tres, cuatro y cinco forma parte del carácter
analítico; no se deduce de la mera existencia del catálogo finito.

Esta precisión elimina una ambigüedad habitual. Una combinación de
invariantes de la construcción es una definición interna y reproducible; su
naturalidad requiere las reglas del lector que especifican orden, signos y
normalización. Aquí esas reglas son visibles y quedan sometidas a los mismos
controles negativos que el resto de la construcción.

\subsection{Fase orientada y coordenada angular conjugada}

La fase de quinto grado anterior al retorno es
\[
y_5=\frac\pi{90}\bigl(H_5+\alpha\bigr).
\]
El carácter regional completo produce la corrección
\(\mathscr C_\pi(\alpha)\). Con la orientación fijada en la
\cref{p12-83-c13-s3:def:lector-regional-determinantal}, definimos
\[
\boxed{
y_*=\frac\pi{90}\bigl(H_5+\alpha\bigr)
-169\alpha^6
-\frac{D_A\alpha^7}{1-D_A\alpha}.}
\]
Ésta es la fórmula primaria. La unidad sexagesimal no interviene hasta aplicar la carta
\[
C_*^{\rm HMT}=\frac{180}{\pi}y_*.
\]

\begin{theorem}[Coordenada angular conjugada del lector regional determinantal]
\label{p12-83-c13-s3:thm:contraangulo-regional}
Fijados el cierre dodecafásico de \(\alpha\),
\(\varphi=(1+\sqrt5)/2\) y las reglas del lector regional determinantal, la
fase \(y_*\) y su coordenada angular conjugada quedan determinados sin introducir el valor
buscado. Su evaluación es
\[
\boxed{
y_*=0.03706622646583826398493779409230638\ldots,}
\]
\[
\boxed{
C_*^{\rm HMT}
=2.12373833896864572426195138039336\ldots{}^\circ.}
\]
Por tanto, a quince cifras decimales,
\[
\boxed{C_*^{\rm HMT}=2.123738338968646^\circ.}
\]
\end{theorem}

\begin{proof}
El \cref{p12-83-c13-s3:thm:unicidad-retorno-169} determina \(169\). El cierre de
\(\alpha\) determina \(A\), \(x\) y
\(D_A=e^{-100\pi\alpha/9}\). El \cref{p12-83-c13-s3:thm:realizacion-analitica-graduada}
determina la serie completa. La fórmula de \(H_5\) contiene únicamente
\(\alpha\), \(\varphi\) e invariantes estructurales declarados. Sustituir
estos datos en la expresión de \(y_*\) da el valor impreso. El valor decimal
de \(C_*^{\rm HMT}\) no aparece en ninguno de los miembros derechos.
\end{proof}

La componente reducida de acción posterior al retorno queda ahora como
salida:
\[
\boxed{
\bar h_{\rm HMT}^{\rm ret}
=\frac{C_*^{\rm HMT}}2-\alpha
=H_5-\frac{90}{\pi}\mathscr C_\pi(\alpha).}
\]
Numéricamente,
\[
\bar h_{\rm HMT}^{\rm ret}
=1.05457181691503906113369058472430\ldots.
\]
La identidad inversa no se ha usado para introducir \(C_*^{\rm HMT}\); se
obtiene después de generar \(y_*\).

\subsection{Truncamiento de grado \(7\) y cota exacta del resto}

La expresión hallada inicialmente,
\[
y_7=\frac\pi{90}(H_5+\alpha)-169\alpha^6-D_A\alpha^7,
\]
es el truncamiento de grado siete del carácter completo. Produce
\[
C_7=2.12373833896864600265403827103919\ldots{}^\circ.
\]
El resto omitido no es un término indeterminado; es la cola geométrica
exacta
\[
\mathscr C_\pi(\alpha)
-169\alpha^6-D_A\alpha^7
=\frac{D_A^2\alpha^8}{1-D_A\alpha}.
\]
Por tanto,
\[
\left|C_7-C_*^{\rm HMT}\right|
=\frac{180}{\pi}\frac{D_A^2\alpha^8}{1-D_A\alpha}
=2.7839208689064583\ldots\times10^{-16}\;{}^\circ.
\]
La recurrencia transforma así una aproximación de séptimo grado en una
realización analítica completa y proporciona una cota cerrada para cada
truncamiento posterior.

\subsection{Ablaciones del selector y de la continuación}

La construcción cambia al retirar cualquiera de sus componentes activos.
Los siguientes valores se obtienen sin recalibrar los demás términos:
\[
\begin{array}{@{}lc@{}}
\toprule
\text{variante}&\text{coordenada angular conjugada en grados}\\
\midrule
\rho_\pi=168&2.1237383389772976863904487\ldots\\
\rho_\pi=169&2.1237383389686457242619514\ldots\\
\rho_\pi=170&2.1237383389599937621334541\ldots\\
\text{sin cola determinantal}&2.1237383389686949415301675\ldots\\
D_A\text{ sustituido por }1&2.1237383389686313409965826\ldots\\
\bottomrule
\end{array}
\]
Estas ablaciones no constituyen por sí solas una estimación probabilística.
Demuestran que el selector equivariante y el carácter determinantal son
componentes efectivas de la salida y que no pueden suprimirse manteniendo la
misma realización.

\subsection{Reconstrucción espectral y transición hexada--octada}

La fase común y la fase orientada determinan los dos canales
\[
q_-=e^{-x+y_*},
\qquad
q_+=e^{-x-y_*}.
\]
Se recuperan exactamente los invariantes
\[
q_-q_+=D_A,
\qquad
\alpha=-\frac9{100\pi}\log D_A,
\qquad
C_*^{\rm HMT}=\frac{90}{\pi}\log\frac{q_-}{q_+}.
\]
La primera igualdad lee la contracción común; la tercera, la separación
orientada de las hojas. El intercambio \(R\mapsto-R\) permuta
\(q_+\) y \(q_-\), conserva \(D_A\) e invierte el signo de \(y_*\).

Para
\[
S_{90}(q)=\frac{q^{90}}{1-q^{270}},
\qquad
S_{120}(q)=\frac{q^{120}}{1-q^{360}},
\]
el funcional de transición hexada--octada se escribe íntegramente en la
carta de fases:
\[
\Gamma_{6\to8}
=\frac{180}{\pi}xy_*
+12\bigl(S_{90}(q_-)-S_{90}(q_+)\bigr)
-\bigl(S_{120}(q_-)-S_{120}(q_+)\bigr).
\]
La evaluación producida por la coordenada \(C_*^{\rm HMT}\) del
\cref{p12-83-c13-s3:thm:contraangulo-regional} es
\[
\boxed{
\Gamma_{6\to8}
=0.274007672694459274747255260774\ldots.}
\]
Así, el valor interno reconocido en Mecánica Dimensional como funcional de
Barbero--Immirzi deja de recibir \(C_*^{\rm HMT}\) como literal independiente:
ambos descienden del mismo par de fases generado en este capítulo. La
interpretación geométrica y física de \(\Gamma_{6\to8}\) pertenece a la obra
de Mecánica Dimensional; su evaluación matemática pertenece ya a la interfaz
previamente construida.

\subsection{Acción anterior y posterior al retorno}

El carácter de quinto grado y la componente de acción posterior al retorno
son dos valores distintos:
\[
H_5
=1.05457181764615446962582182943306\ldots,
\]
\[
\bar h_{\rm HMT}^{\rm ret}
=1.05457181691503906113369058472430\ldots.
\]
El primero es la acción local antes de incorporar el carácter regional; el
segundo es la acción orientada que queda después del retorno. La corrección
que genera la coordenada angular conjugada impide identificarlos.

La comparación metrológica debe respetar esta distinción. El primer valor da
la mantisa
\[
2\pi H_5
=6.62607014999998792600111034989947\ldots,
\]
mientras que la acción posterior da
\[
2\pi\bar h_{\rm HMT}^{\rm ret}
=6.62607014540625433351074984759288\ldots.
\]
En el Sistema Internacional, el valor
\[
h=6.62607015\times10^{-34}\;\mathrm{J\,s}
\]
es exacto por definición \cite{BIPMSI}. Por ello,
\[
2\pi H_5-6.62607015
=-1.2073998889650101\ldots\times10^{-14},
\]
y no es una igualdad exacta. La proximidad del valor anterior al retorno es
un control numérico notable; no autoriza a sustituir la constante definitoria
del SI por la componente posterior ni a confundir las dos acciones HMT\@.

Este resultado determina con precisión la situación matemática. El lector
regional genera \(C_*^{\rm HMT}\) y, desde él, el funcional
\(\Gamma_{6\to8}\). El mismo cálculo separa la aproximación de Planck
producida por \(H_5\) de la acción reducida que interviene en el
la coordenada angular conjugada. La separación no debilita la construcción: elimina una
identificación incompatible con las cifras exactas y convierte cada
comparación en una afirmación falsable de tipo bien definido.

\subsection{Transportador relativo y disjunción espectral de los bloques}
\label{p12-83-c13-s3:sec:transportador-relativo-split}

En la base ordenada de canales, sea
\[
 R=\begin{pmatrix}0&1\\1&0\end{pmatrix},
 \qquad P_\pm=\frac{I\pm R}{2},
 \qquad
 B_0=7I+2R,
 \qquad B_1=\Adeg I+\Cdeg R.
\]
La descomposición escindida
\[
 xI+yR=\rho(x,y)e^{\eta(x,y)R},
 \quad \rho(x,y)=\sqrt{x^2-y^2},
 \quad \eta(x,y)=\operatorname{artanh}(y/x)
\]
construye el transportador único
\[
 T_{\rm rel}
 =\frac{\sqrt{(\Adeg)^2-(\Cdeg)^2}}{\sqrt{45}}
   \exp\!\left[
   \left(\operatorname{artanh}\frac{\Cdeg}{\Adeg}
   -\operatorname{artanh}\frac27\right)R\right],
 \qquad T_{\rm rel}B_0=B_1.
\]
En la base espectral actúa por
\(9\mapsto\Adeg+\Cdeg\) y
\(5\mapsto\Adeg-\Cdeg\).

\begin{theorem}[Disjunción espectral de las realizaciones local y angular]
\label{p12-83-c13-s3:thm:no-go-entrelazador-bloques-split}
Los espectros
\[
 \operatorname{Spec}(B_0)=\{9,5\},\qquad
 \operatorname{Spec}(B_1)=\{\Adeg+\Cdeg,\Adeg-\Cdeg\}
\]
son disjuntos. En consecuencia,
\[
 \operatorname{Hom}_{B_0,B_1}
 :=\{\Phi:\Phi B_0=B_1\Phi\}=\{0\}.
\]
\end{theorem}

\begin{proof}
Las cotas racionales construidas para las coordenadas angulares dan
\(7.29<\Adeg<7.30\) y \(2.12<\Cdeg<2.13\); por tanto,
\(9.41<\Adeg+\Cdeg<9.43\) y
\(5.16<\Adeg-\Cdeg<5.18\). En la base que diagonaliza \(R\), cada entrada
de un entrelazador queda multiplicada por la diferencia entre un autovalor
de \(B_0\) y uno de \(B_1\); la disjunción obliga a que todas se anulen.
\end{proof}

El transportador relativo compara dos bloques dentro del álgebra escindida;
los entrelazadores dodecafásico--Witt y de incidencia actúan entre
representaciones equivalentes en sus residencias propias. Esta separación
tipa positivamente ambas operaciones y conserva sus dominios.

\subsection{Cadena de generación \(\mathcal F_\pi\to b_\pi\to\rho_\pi\to D_A\to C_*^{\rm HMT}\to\Gamma_{6\to8}\): selección regional, prolongación determinantal y coordenada orientada}

El procedimiento completo puede resumirse como una cadena de aplicaciones:
\[
\begin{aligned}
\mathcal F_\pi
&\longmapsto b_\pi
\longmapsto\rho_\pi=169,\\
\alpha
&\longmapsto A
\longmapsto x
\longmapsto D_A,\\
(\rho_\pi,D_A)
&\longmapsto\mathscr C_\pi(\alpha),\\
(H_5,\mathscr C_\pi(\alpha))
&\longmapsto y_*
\longmapsto C_*^{\rm HMT},\\
(x,y_*)
&\longmapsto(q_+,q_-)
\longmapsto\Gamma_{6\to8}.
\end{aligned}
\]
La primera línea pertenece al catálogo finito; la segunda, al cierre común de
\(\alpha\); la tercera, a la prolongación determinantal; la cuarta produce la
coordenada orientada; la quinta la transporta hacia la interfaz dimensional.

En esta organización, \(\alpha\) mide la contracción común del transporte y
\(C_*^{\rm HMT}\) mide la separación orientada de sus dos hojas. El retorno
regional no añade una constante externa: selecciona una amplitud finita y la
prolonga mediante el determinante ya fijado por \(\alpha\). La carta en
grados es posterior. El objeto primario es el par de fases \((x,y_*)\).

Una verificación exhaustiva reconstruye las cinco regiones directamente
desde el catálogo, recorre las acciones de
\(\mathfrak S_5\) y \(\mathfrak S_3\), comprueba la recurrencia y suma la
serie. La entrada de \(\alpha\) es la salida exacta de su cierre
dodecafásico, no una cifra metrológica transcrita. La comprobación mantiene
separados los teoremas del catálogo, las reglas del lector y el contraste
externo; por ello ninguna comparación experimental interviene en la función
generadora.
